%% P07 --- Tensory, kształty i notacja wskaźnikowa
%%
%% Każda liczba w tym programie, której czytelnik nie policzy w pamięci,
%% pochodzi z code/p07_tensors_shapes.py i trafia na stronę przez \val{}.
%% Listing konsolowy to zatwierdzony transkrypt pisany przez ten skrypt.
%%
%% TEZA: działanie na tablicach jest stwierdzeniem o tym, które osie się
%% zestawiają, a notacja wskaźnikowa mówi to dokładnie tam, gdzie obrazek nie
%% potrafi.
%%
%% DLACZEGO TEN PROGRAM ISTNIEJE. Został dodany przez przegląd programu
%% nauczania z sierpnia 2026, a nie zaprojektowany, i przegląd nazwał go
%% największą luką treściową w tej książce i w każdej z przeglądu
%% konkurencyjnego: odbiorca codziennie operuje tablicami rzędu 4, uczono go
%% wyłącznie notacji rzędu 2, a z tej rozbieżności biorą się błędy kształtu.
%% Plan ramka po ramce to notes/08-tensors-plan.md, napisany przed pierwszym
%% wierszem tego pliku.
%%
%% CO TEN PROGRAM JEST WINIEN, odczytane z plików, a nie z pamięci:
%%   P06  Sekcja 4 to cała historia rzędu 2 i przekazuje temu programowi
%%        cztery rzeczy po imieniu. Kształt to zapisana dziedzina i
%%        przeciwdziedzina; błąd kształtu jest więc błędem TYPU; partia k
%%        wejść ma kształt n x k w konwencji P06 -- a jego nota mówi, że
%%        prawdziwe frameworki układają partię wzdłuż PIERWSZEJ osi, więc
%%        warstwa liczy X W^T z X o kształcie k x n, i że to P07 nazywa osie.
%%        Ostatnia ramka P06 przekazuje pytanie o cztery indeksy.
%%        TEN PROGRAM ZACZYNA WIĘC OD ZAMIANY KONWENCJI.
%%   P04  ma wymiar jako liczbę niezależnych kierunków, czyli słowo, z którym
%%        ten program się zderza. Nazwij to zderzenie.
%%   P03  ma rachunek pamięci i intensywność arytmetyczną, więc koszt kopii
%%        jest stamtąd cytowany i nigdy tu mierzony ponownie.
%%
%% CZEGO TEN PROGRAM NIE RUSZA, sprawdzone w tools/programs.json:
%%   rząd, cztery podprzestrzenie, najmniejsze kwadraty, LoRA         -> P08
%%   wyznacznik, odwrotność, zmiana bazy                              -> P09
%%   transformer złożony z tego wszystkiego                           -> P32
%%   różniczkowanie przez reshape                                     -> P16, P18
%%
%% Bliźniak: programs/en/P07-tensors-shapes.tex. Ramka za ramką, makro za
%% makrem, liczba za liczbą; tools/parity.py je porównuje.

\program{Tensory, kształty i notacja wskaźnikowa}
\label{prog:P07}
\index{tensor!shape}
\index{tensor!index notation}

\begin{outcomes}
\outcome{1--8}{czytać kształt jako krotkę długości osi i mówić, co liczy dana
  wyjęta z niego liczba}
\outcome{9--15}{zapisać działanie na tablicach w notacji wskaźnikowej i z
  samych wskaźników powiedzieć, które osie je przeżywają}
\outcome{16--22}{czytać napis \code{einsum} jako sumę, którą oznacza, i pisać
  taki napis dla działania podanego słowami}
\outcome{23--30}{stosować regułę rozgłaszania dokładnie i rozpoznawać
  kształty, które zestawiają się w sposób, o który nikt nie prosił \dash{} ta
  sekcja zapisuje swoją tożsamość przez $\operatorname{Cov}(p, t)$ i
  $\operatorname{Var}(t)$, które definiuje Program~\ref{prog:P24}, i podaje
  jedno zdanie o każdym z nich, którego wymaga argument}
\outcome{31--37}{powiedzieć, które z reshape, transpozycji i permute przesuwa
  dane, i która kolejność poprawnie dzieli partię na głowice}
\end{outcomes}

\begin{quiz}
\item Tablica ma kształt $(\val{p07.emb.b}, \val{p07.emb.s}, \val{p07.emb.d})$.
  Ile to liczb i ile osi? \teachesat{7--8}
  \answerto{$\val{p07.emb.total}$ liczb na $\val{p07.emb.rank}$ osiach. Liczba
  wpisów to iloczyn długości; liczba osi to ile jest tych długości, a te dwa
  pytania są różne.}
\item W $\sum_k A_{ik} B_{kj}$ które litery lądują na wyniku?
  \teachesat{9--10}
  \answerto{$i$ oraz $j$. Wskaźnik, który pojawia się dwa razy, jest sumowany
  i znika; wskaźnik, który pojawia się raz, przeżywa. Ten jeden wiersz to cała
  notacja wskaźnikowa.}
\item Co robi napis \code{einsum} postaci \code{'ij->ji'}? \teachesat{17--18}
  \answerto{Transponuje. Oba wskaźniki przeżywają i żaden nie jest sumowany;
  zmieniła się tylko kolejność, w której są zapisane.}
\item Dwa kształty, $(n)$ i $(n, 1)$. Jaki kształt daje odjęcie jednego od
  drugiego? \teachesat{25--26}
  \answerto{$(n, n)$. Zestawienie kształtów od prawej łączy $n$ z $1$, co się
  rozciąga, a wiodące $n$ nie ma z czym się połączyć, więc zostaje.
  Powstają wszystkie pary i nic nie protestuje.}
\item Które z reshape, transpozycji i permute przesuwa liczby w pamięci?
  \teachesat{31--32}
  \answerto{Żadne, w typowym przypadku. Każde oddaje ten sam ciąg liczb z inną
  regułą czytania. Dane przesuwa dopiero żądanie, żeby wynik był spójny w
  pamięci.}
\item Tablica o kształcie $(b, s, h \times d)$ ma stać się $(b, h, s, d)$. Czy
  wystarczy jeden reshape? \teachesat{33--34}
  \answerto{Nie. Reshape musi rozbić ostatnią oś, co daje $(b, s, h, d)$, a
  dwie środkowe osie trzeba potem zamienić. Reshape wprost do $(b, h, s, d)$
  jest legalny, cichy i daje inną tablicę.}
\end{quiz}

% ============================================================
\section{Tablica to kształt i reguła czytania}
\label{sec:P07-shape}
\index{tensor!shape}

\begin{fr}
Program~\ref{prog:P06} zostawił jedną rzecz niedokończoną, i to celowo.

Pokazał, że partia $k$ wejść przez warstwę $W$ jest macierzą o $k$ kolumnach,
po jednym wejściu na kolumnę, a potem zauważył, że prawdziwe frameworki tak nie
robią: układają partię wzdłuż \emph{pierwszej} osi i liczą $XW\T$ z $X$ o
kształcie $k \times n$.

Dwie konwencje, a tylko jedna z nich może być tą warstwą. Która?
\dotline
\nextframe
\end{fr}

\begin{fr}
\ans{Obie.}
$XW\T$ jest transpozycją $WX$, a transpozycja nie jest innym rachunkiem. Obie
konwencje opisują jedno przekształcenie i różnią się tylko tym, którą stroną
do góry je zapisać.

Warto mieć to jako pierwszą ramkę programu, bo to ostatni raz, kiedy pytanie o
tablice da się rozstrzygnąć tożsamością macierzową.
Program~\ref{prog:P06} miał dwa indeksy i każde jego pytanie mieściło się w
nich. Tensor uwagi ma cztery.
\end{fr}

\begin{fr}
Pierwszym zadaniem jest więc definicja, która przeżyje cztery.

\emph{Kształt} tablicy to krotka długości. Kształt $(2, 3)$ to dwa wiersze po
trzy, a $(\val{p07.emb.b}, \val{p07.emb.s}, \val{p07.emb.d})$ to partia zdań
złożonych z wektorów. \emph{Oś} jest pozycją w tej krotce.

\textbf{Oś jest pozycją w krotce indeksów, a nie kierunkiem w przestrzeni.}
Wszystko inne w tym programie to zastosowanie tego zdania.
\end{fr}

\begin{fr}
I tu jedno słowo robi trzy rzeczy naraz, które się nie zgadzają.

Tablica o kształcie $(\val{p07.emb.b}, \val{p07.emb.s}, \val{p07.emb.d})$ to
partia $\val{p07.emb.b}$ zdań, każde po $\val{p07.emb.s}$ tokenów, każdy z
jednym embeddingiem.

Konkretnie: pierwsza oś biegnie po zdaniach w partii, druga po pozycjach tokenów
w zdaniu, a trzecia po liczbach tworzących embedding jednego tokenu. Trzy osie i
trzy całkiem różne rzeczy liczone wzdłuż nich.

Embeddingi są wektorami. Program~\ref{prog:P04} zapytałby o wymiar przestrzeni,
w której one żyją. Powiedz, ile on wynosi.
\dotline
\nextframe
\end{fr}

\begin{fr}
\ans{$\val{p07.emb.d}$.}
Wektory mają po $\val{p07.emb.d}$ współrzędnych, żyją więc w przestrzeni o
$\val{p07.emb.d}$ wymiarach i stosuje się do nich cały aparat
Programu~\ref{prog:P04}.

\begin{trapbox}
\textbf{Nazwanie tego tablicą trójwymiarową i rozumowanie potem o trzech
wymiarach jest błędem, a szkodę robi słowo.}

Trzy rzeczy w tamtym zdaniu nazywają się wymiarem, a odpowiedź na pytanie
\enquote{jaki jest wymiar} jest dla każdej inna:

\begin{itemize}
\item \textbf{liczba osi} wynosi $\val{p07.emb.rank}$;
\item \textbf{długość osi} to $\val{p07.emb.b}$,
  $\val{p07.emb.s}$ albo $\val{p07.emb.d}$, zależnie która;
\item \textbf{wymiar przestrzeni}, w której żyje zapisany wektor, wynosi
  $\val{p07.emb.d}$ i jest faktem o jednej osi, a nie o tablicy.
\end{itemize}

Obrazek tylko pogarsza sprawę, bo obrazek trzech osi jest obrazkiem trzech
kierunków, a to nie jest to.
\end{trapbox}
\end{fr}

\begin{fr}
\begin{note}
A liczba osi nazywa się \emph{rzędem} tablicy, co zderza się z drugą rzeczą
nazywaną rzędem dwa programy dalej.

Program~\ref{prog:P08} używa tego słowa na liczbę niezależnych kierunków, na
które macierz odwzorowuje, czyli na stwierdzenie o tym, co macierz
\emph{robi}. Rząd tablicy jest stwierdzeniem o tym, jak głęboko sięga
indeksowanie, i o niczym więcej. Tablica rzędu $\val{p07.emb.rank}$ o tamtym
kształcie nie mówi nic o żadnym z nich i żadne odczytanie nie implikuje
drugiego.

Te dwa są dość daleko, żeby zapomnieć, i dość blisko, żeby się zderzyć. Ta
książka pisze \emph{liczba osi} wszędzie tam, gdzie jest jakakolwiek
wątpliwość.
\end{note}
\end{fr}

\begin{fr}
Pod kształtem leży jeden płaski ciąg liczb, a to znaczy więcej, niż brzmi.

Tablica o kształcie $(\val{p07.emb.b}, \val{p07.emb.s}, \val{p07.emb.d})$
mieści ile liczb?
\dotline
\nextframe
\end{fr}

\begin{fr}
\ans{$\val{p07.emb.total}$.}
Iloczyn długości, czyli reguła liczenia, którą dał ci Program~\ref{prog:F10}.

Leżą w pamięci w jednej linii, a kształt jest regułą zamieniającą krotkę
indeksów na pozycję w tej linii. Dla kształtu $(n_0, n_1, n_2)$ wpis o
indeksach $(i, j, k)$ leży pod

\[ i \, n_1 n_2 + j \, n_2 + k \]

więc ostatni indeks przesuwa się o jeden krok naraz, a pierwszy skacze co
$n_1 n_2$.

\textbf{Tablica to bufor, kształt i ta reguła.} Sekcja 5 pyta o to, które z
trzech działań zmienia które z tych trzech, i bez tej ramki nie da się jej
nawet postawić.
\end{fr}

\mermaidfig{p07-axes-not-directions}{Czym jest oś i która z trzech rzeczy
zwanych wymiarem jest która.}{osie to nie kierunki}

% ============================================================
\section{Notacja wskaźnikowa mówi to, czego obrazek nie potrafi}
\label{sec:P07-indices}
\index{tensor!index notation}

\begin{fr}
Program~\ref{prog:P06} wyprowadził iloczyn macierzy jako złożenie. Zapisany we
wskaźnikach jest to

\[ C_{ij} = \sum_k A_{ik} B_{kj} \]

czyli to samo stwierdzenie i inny akt: złożenie mówi \emph{dlaczego}, a to mówi
\emph{które}.

Patrz na litery, a nie na sumę. $i$ oraz $j$ pojawiają się po prawej raz każde;
$k$ pojawia się dwa razy. Co robi ta różnica?
\dotline
\nextframe
\end{fr}

\begin{fr}
\begin{ansblock}
Wskaźnik, który pojawia się dwa razy, jest sumowany i znika; wskaźnik, który
pojawia się raz, przeżywa na wynik.
\end{ansblock}

To cała reguła i warto powiedzieć ją płasko, bo wszystko, co zostało w tym
programie, jest nią zastosowaną do większej liczby liter.

Sigma jest wtedy prawie ozdobą: skoro $k$ pojawia się dwa razy, wiesz, że jest
sumowane, a skoro $i$ oraz $j$ pojawiają się raz, wiesz, że wynik jest
indeksowany przez $i$ oraz $j$. \textbf{Kształt wyniku da się odczytać z
samych liter, bez żadnej arytmetyki.}
\end{fr}

\begin{fr}
Przeczytaj więc jedno z powrotem, a to nie jest iloczyn.

\[ D_{ij} = \sum_k A_{ik} B_{jk} \]

Dwa wskaźniki $B$ zamieniły się miejscami. Czym jest $D$?
\dotline
\nextframe
\end{fr}

\begin{fr}
\ans{$AB\T$.}
$B_{jk}$ to wpis w wierszu $j$ i kolumnie $k$, więc sumowanie po drugim
wskaźniku obu zestawia wiersz $i$ macierzy $A$ z wierszem $j$ macierzy $B$
\dash{} czyli $A$ razy $B$ transponowane.

Co zamyka pętlę otwartą na początku programu: $XW\T$ jest warstwą frameworka, a
we wskaźnikach jest to $\sum_k X_{ik} W_{jk}$.

Ile więc danych przenosi ta transpozycja?
\dotline
\nextframe
\end{fr}

\begin{fr}
\ans{Żadnych. \textbf{Transpozycja nigdy nie musiała się zdarzyć.} Jest opisem
tego, który wskaźnik spotyka który, i żadna tablica nie musi być obracana, żeby
go spełnić.}
Teraz cztery litery, na obiekcie, dla którego ten program istnieje.

\[ S_{bhqk} = \sum_d Q_{bhqd} K_{bhkd} \]

Nieważne, jak to się nazywa. Zastosuj regułę: powiedz, które wskaźniki
przeżywają na $S$, a które są wysumowane.
\dotline
\nextframe
\end{fr}

\begin{fr}
\begin{ansblock}
$b$, $h$, $q$ i $k$ przeżywają; sumowane jest tylko $d$.
\end{ansblock}

$d$ jest jedyną literą, która pojawia się dwa razy \emph{i} nie ma jej po
lewej, jest więc jedynym zwężeniem. Wynik ma cztery osie, a jego kształt to
cztery długości $b$, $h$, $q$ i $k$ w kolejności, w jakiej zostały zapisane.

\textbf{Właśnie odczytałeś wsadowy wielogłowicowy wynik uwagi z samych liter,
nie wiedząc, co która oznacza.} To argument za tą notacją w jednej ramce:
obrazek tego obiektu nie istnieje, a zdanie, które go opisuje, mieści się w
jednej linii.
\end{fr}

\begin{fr}
Jedno doprecyzowanie, bo $b$ i $h$ też pojawiają się dwa razy i nie są sumowane.

Reguła w tej postaci jest zbyt zgrubna. Rozstrzyga lewa strona: wskaźnik
nieobecny po lewej jest sumowany, a wskaźnik obecny po lewej przeżywa,
niezależnie od tego, jak często pojawia się po prawej. $b$ i $h$ są po lewej,
są więc wskaźnikami \emph{wsadowymi} \dash{} ten sam rachunek wykonany raz dla
każdego z nich.

Usuń $b$ z lewej strony, a suma pobiegnie i po nim, dając tablicę o trzech
osiach zamiast czterech i dodając do siebie rzeczy należące do różnych
przykładów. Nic nie ostrzega. Sekcja 3 jest miejscem, gdzie to staje się czymś,
co da się wpisać.
\end{fr}

% ============================================================
\section{Napis \texorpdfstring{\code{einsum}}{einsum} to tamto zdanie z pominiętą sigmą}
\label{sec:P07-einsum}
\index{tensor!einsum@\texttt{einsum}}

\begin{fr}
Oto całość.

\begin{center}
\code{'ik,kj->ij'}
\end{center}

Przecinek rozdziela argumenty, strzałka wprowadza wynik, a każda grupa liter to
wskaźniki jednej tablicy po kolei. Porównaj to z $\sum_k A_{ik} B_{kj}$ i nie
ma w jednym niczego, czego nie ma w drugim.

Strzałka wykonuje jedną robotę i jest to ta robota, której potrzebowało
doprecyzowanie z poprzedniej ramki. Powiedz, co napis robi z literą, która
pojawia się wśród argumentów, ale nie za strzałką.
\dotline
\nextframe
\end{fr}

\begin{fr}
\ans{Sumuje ją. \textbf{Wszystko na lewo od strzałki, czego nie ma na prawo,
jest sumowane} \dash{} konwencja zapisana jako składnia, i po to właśnie jest
ta strzałka.}
Przeczytaj więc trzy na zimno. Powiedz, co robi każdy z nich, przy tablicach
dwuwymiarowych tam, gdzie potrzebne są dwie.

\begin{center}
\code{'ij->ji'} \qquad \code{'ii->'} \qquad \code{'ij,j->i'}
\end{center}
\dotline
\nextframe
\end{fr}

\begin{fr}
\begin{ansblock}
Transpozycja; ślad; macierz razy wektor.
\end{ansblock}

Pierwszy zachowuje oba wskaźniki i zmienia ich kolejność. Drugi ma $i$ dwa razy
i nic po strzałce, więc sumuje po $i$ przy związanych ze sobą wskaźnikach obu
argumentów \dash{} przekątna, dodana do siebie. Trzeci zestawia drugi wskaźnik
macierzy z jedynym wskaźnikiem wektora i zachowuje pierwszy.

\textbf{Żaden z tych trzech nie potrzebował nazwy, żeby dał się przeczytać.}
Zastosowałeś jedną regułę trzy razy, a trzeci to zdanie
Programu~\ref{prog:P06} \enquote{wektor jest macierzą o jednej kolumnie},
przybyłe bez tej kolumny.
\end{fr}

\begin{fr}
Teraz w drugą stronę, i to tę wartą umieć.

Masz $Q$ i $K$, oba o kształcie (partia, głowica, pozycja, cecha), i chcesz dla
każdej partii i każdej głowicy iloczyn skalarny każdej pozycji zapytania z
każdą pozycją klucza. Napisz ten napis.
\dotline
\nextframe
\end{fr}

\begin{fr}
\ans{\code{'bhqd,bhkd->bhqk'}}
Pozycje zapytań to $q$, pozycje kluczy to $k$, oś cech $d$ jest tą sumowaną, bo
iloczyn skalarny sumuje po cechach, a $b$ i $h$ są po obu stronach, bo całość
wykonuje się raz na element partii i raz na głowicę.

To suma z ramki 12, wpisana.

Powiedz więc, czym jest ten napis w kategoriach matematyki, którą już masz.
\dotline
\nextframe
\end{fr}

\begin{fr}
\ans{Sigmą z pominiętą sigmą. \textbf{Nie jest to nowy pomysł, sztuczka
biblioteki ani udogodnienie wydajnościowe} \dash{} kto umie przeczytać tę sumę,
przeczyta i ten napis.}
Zostaje ten, który ludzie mylą, i warto w niego wejść.

Rozważ \code{'bij,bjk->bik'}. Trzy litery pojawiają się dwa razy: $b$, $j$ i
jeszcze jedna. Które z nich są sumowane?
\dotline
\nextframe
\end{fr}

\begin{fr}
\ans{Tylko $j$.}
$b$ pojawia się dwa razy po lewej \emph{i} po prawej, jest więc wskaźnikiem
wsadowym: ten sam iloczyn macierzy wykonany raz na element partii. $i$ oraz $k$
pojawiają się po razie. To jest więc partia iloczynów macierzy, czyli to, do
czego ta notacja służy częściej niż do czegokolwiek innego.

\begin{trapbox}
\textbf{\enquote{Powtórzony znaczy sumowany} to reguła, którą ludzie stąd
wynoszą, i jest to reguła dla \emph{sumy}, a nie dla napisu.}

Bierze się z prawdziwej konwencji \dash{} w tradycyjnej konwencji sumacyjnej
nie ma strzałki, więc powtórzony rzeczywiście znaczy sumowany i nic innego
znaczyć nie może. Strzałka jest dokładnie tym, co czyni powtórzony wskaźnik
opcjonalnym, a opuszczenie litery po prawej stronie jest sposobem na
poproszenie o sumę.

Koszt pomylenia tego nie jest błędem. \code{'bij,bjk->ik'} działa, zwraca
tablicę o kształcie właściwym dla pojedynczego przykładu i dodało do siebie
wyniki należące do różnych elementów partii. Strata spada. Model się nie uczy.
\end{trapbox}
\end{fr}

% ============================================================
\section{Rozgłaszanie jest regułą, a nie uprzejmością}
\label{sec:P07-broadcast}
\index{tensor!broadcasting}

\begin{fr}
Dodanie wektora $(3)$ do macierzy $(3, 3)$ nie jest zdefiniowane przez nic, co
spotkałeś, a każda biblioteka tablicowa robi to mimo wszystko. To, co robi,
nazywa się \emph{rozgłaszaniem} i ma dokładną regułę.

\begin{center}
Zestaw oba kształty \textbf{od prawej}. Uzupełnij krótszy z przodu wartościami
$1$. Para osi zgadza się, gdy długości są równe albo gdy jedna z nich wynosi
dokładnie $1$; jeśli któraś para nie spełnia żadnego z tych warunków, działanie
zostaje odrzucone.
\end{center}

Trzy zdania, a rozstrzygają każde pytanie o kształt w tej sekcji. Przeczytaj je
jeszcze raz, zamiast parafrazować \dash{} parafraza, którą się wynosi, brzmi
\enquote{robi rzecz rozsądną}, a tego nie robi.
\end{fr}

\mermaidfig{p07-broadcast-alignment}{Jak zestawia się dwa kształty i co liczy
się jako zgodność.}{zestawianie kształtów}

\begin{fr}
Zatem $(3)$ przeciwko $(3, 3)$: uzupełnione do $(1, 3)$, pary to $1$ z $3$, co
się rozciąga, oraz $3$ z $3$, co pasuje. Wektor zostaje dodany wzdłuż każdego
wiersza.

Załóż, że chcesz go dodać w dół każdej kolumny. Te same trzy liczby, ta sama
macierz i żadna arytmetyka się nie zmienia. Co musi się zmienić?
\dotline
\nextframe
\end{fr}

\begin{fr}
\ans{Kształt, który nadajesz wektorowi: $(3, 1)$.}
Pary stają się $3$ z $3$ oraz $1$ z $3$, więc rozciąganie dzieje się w drugą
stronę.

Oba są legalne, oba działają i tylko jedno jest tym, o co ci chodziło. Nic w
arytmetyce nie powie ci które. \textbf{Kształt jest całą instrukcją, a $(3)$ i
$(3, 1)$ to różne instrukcje.}

Co przygotowuje ten, który kosztuje pieniądze, i warto odpowiedzieć, zanim
pójdziesz dalej.

Masz predykcje w tablicy o kształcie $(n)$ i wartości docelowe w tablicy o
kształcie $(n, 1)$ \dash{} rzecz zupełnie zwyczajna, bo jedno wyszło z modelu,
a drugie z kolumny tabeli. Odejmujesz jedno od drugiego.

Zastosuj regułę. Jaki kształt wychodzi?
\dotline
\nextframe
\end{fr}

\begin{fr}
\ans{$(n, n)$.}
Uzupełnione do $(1, n)$, pary to $1$ przeciwko $n$ oraz $n$ przeciwko $1$. Obie
się rozciągają. Zamiast $n$ różnic masz więc ich $n^{2}$ \dash{} każdą predykcję
przeciwko każdej wartości docelowej.

Potem podnosisz do kwadratu i bierzesz średnią, a średnia z $n^{2}$ liczb jest
liczbą, więc strata jest liczbą, więc pętla treningowa jest zupełnie
zadowolona.

\textbf{Nie ma błędu, nie ma ostrzeżenia i nie ma kształtu do obejrzenia.} Dwie
następne ramki są o tym, jak bardzo ta liczba jest zła, a odpowiedź nie brzmi
\enquote{trochę}.
\end{fr}

\begin{fr}
Da się to policzyć dokładnie, a nie zmierzyć, co dla błędu jest rzadkie.

Pisząc $p$ na predykcje i $t$ na wartości docelowe, średnia po wszystkich
parach rozwija się w trzy składniki, a wychodzi z tego

\[ \frac{1}{n^{2}}\sum_{i,j} (p_i - t_j)^{2}
   \;=\; \frac{1}{n}\sum_{i} (p_i - t_i)^{2} \;+\; 2\operatorname{Cov}(p, t). \]

Raportowana strata to strata prawdziwa plus podwojona kowariancja predykcji z
wartościami docelowymi.

Teraz pytanie, dla którego ta tożsamość jest. Załóż, że model staje się
doskonały, więc każda predykcja równa się swojej wartości docelowej. Co robi
raportowana strata?
\dotline
\nextframe
\end{fr}

\begin{fr}
\ans{Zatrzymuje się na $2\operatorname{Var}(t)$.}
Przy doskonałym dopasowaniu prawdziwa strata wynosi zero, a kowariancja jest
wariancją wartości docelowych, więc raportowana liczba wynosi
$2\operatorname{Var}(t)$ i nie zejdzie niżej, jakkolwiek dobry byłby model.

\transcript{p07-broadcast}

Trzy predykcje równe swoim trzem wartościom docelowym: prawdziwa strata wynosi
$0$, a raportowana $\val{p07.demo.floor}$, czyli podwojoną wariancję
$1, 2, 3$.

\textbf{Inżynier widzi stratę, która spada, wypłaszcza się na liczbie, której
nikt nie umie wyjaśnić, i tam zostaje.} Czyta się to jako model, który
przestał się uczyć, a jest to brakujący argument w jednym wierszu.
\end{fr}

\begin{fr}
Zmierzone na $\val{p07.mse.n}$ przykładach przy dwóch poziomach szumu, a
kierunek jest tym, co warto zapamiętać.

\begin{center}
\begin{tabular}{lrr}
\toprule
 & strata prawdziwa & raportowana \\
\midrule
słabe dopasowanie & $\val{p07.mse.true.hi}$ & $\val{p07.mse.bcast.hi}$ \\
dobre dopasowanie & $\val{p07.mse.true.lo}$ & $\val{p07.mse.bcast.lo}$ \\
\bottomrule
\end{tabular}
\end{center}

Przy słabym dopasowaniu raportowana strata jest $\val{p07.mse.ratio.hi}$ razy
większa od prawdziwej. Przy dobrym jest $\val{p07.mse.ratio.lo}$ razy większa.

Zestaw te dwa wiersze ze sobą i powiedz, w którą stronę biegnie błąd, w miarę
jak model się poprawia, i dlaczego.
\dotline
\nextframe
\end{fr}

\begin{fr}
\ans{Robi się gorzej. \textbf{Im lepszy staje się model, tym bardziej
raportowana liczba kłamie}, bo nadwyżka jest podwojoną kowariancją, a
kowariancja jest dokładnie tym, co trening zwiększa.}
Błąd łagodny, dopóki debugujesz, i ostry, gdy wszystko już działa, ma najgorszy
możliwy kształt.

\begin{aibox}
Poprawka to jeden argument i warto wiedzieć, co robi, a nie tylko że działa.
Redukcja tablicy wzdłuż osi normalnie tę oś usuwa: $(n, 1)$ zsumowane wzdłuż
drugiej osi daje $(n)$. Prośba o jej zachowanie zostawia $(n, 1)$, a różnica
jest dokładnie tą, o której jest ta sekcja.

\code{keepdims} nie jest więc flagą formatującą. Jest decyzją o zostawieniu osi
na miejscu, żeby zestawienie w następnym działaniu było wypowiedziane, a nie
wywnioskowane \dash{} a powodem, żeby sięgać po nią domyślnie, jest to, że
zachowana oś długości $1$ może się tylko rozciągnąć, podczas gdy oś brakująca
zestawia się z tym, co akurat naprzeciw.
\end{aibox}
\end{fr}

% ============================================================
\section{Które z nich przesuwa dane}
\label{sec:P07-reshape}
\index{tensor!reshape}

\begin{fr}
Trzy działania zmieniają kształt i nie są tym samym działaniem.

\emph{Reshape} zostawia liczby w kolejności, w jakiej są, i czyta je pod nowym
kształtem. \emph{Transpozycja} zamienia dwie osie. \emph{Permute} zmienia
kolejność wszystkich naraz, a transpozycja jest przypadkiem szczególnym z
dwiema.

Sekcja 1 powiedziała, że tablica to bufor, kształt i reguła zamieniająca krotkę
indeksów na pozycję. Które z tych trzech przesuwa wobec tego liczby w pamięci?
\dotline
\nextframe
\end{fr}

\begin{fr}
\ans{W typowym przypadku żadne.}
Reshape zostawia bufor i zmienia kształt, co sekcja 1 uczyniła oczywistym.
Niespodzianką jest to drugie: transponowanie też niczego nie przesuwa. Zmienia
\emph{regułę}, tak że krotka indeksów zamienia się na pozycję w drugą stronę, i
oddaje ten sam bufor.

Dane przesuwa dopiero prośba o wynik \emph{spójny w pamięci} \dash{} o świeży
bufor, którego kolejność pasuje do kształtu \dash{} a to jest trzecie działanie
z własną nazwą. Program~\ref{prog:P03} wycenił przebieg przez pamięć; tutaj
dowiadujesz się, że o niego poprosiłeś.

\textbf{Zdanie \enquote{reshape jest tani, a transpozycja droga} jest więc
błędne podwójnie}, a wierzy się w nie dlatego, że po transpozycji czasem
następuje kopia, której nikt nie zapisał.
\end{fr}

\begin{fr}
Teraz miejsce, które kosztuje ludzi dzień, i jest to dokładnie to, w którym
partia staje się głowicami.

Masz tablicę o kształcie $(b, s, h \times d)$: partia, pozycja i jeden długi
wektor na pozycję, który ma być $h$ głowicami po $d$ cech. Ma stać się
$(b, h, s, d)$.

Oto dwie drogi do tego. Obie dają ten kształt.

\begin{center}
\code{x.reshape(b, s, h, d).transpose(1, 2)} \\[2pt]
\code{x.reshape(b, h, s, d)}
\end{center}

Czy dają tę samą tablicę?
\dotline
\nextframe
\end{fr}

\begin{fr}
\ans{Nie.}
Na tablicy o $\val{p07.split.entries}$ wpisach przy
$b = \val{p07.split.b}$, $s = \val{p07.split.s}$, $h = \val{p07.split.h}$ i
$d = \val{p07.split.d}$ oba wyniki różnią się na
$\val{p07.split.differ}$ z nich.

Poprawna jest tylko pierwsza. Bufor biegnie najszybciej po $d$, potem po $h$,
potem po $s$, potem po $b$ \dash{} tak mówi kształt $(b, s, h \times d)$ \dash{}
więc rozbicie ostatniej osi na $h$ i $d$ jest przeetykietowaniem tego, co już
tam jest, a dwie środkowe osie trzeba potem zamienić.

Druga prosi o liczby w kolejności, której bufor nie ma, i dostaje ją mimo to,
bo reshape nie sprawdza niczego poza sumą wpisów.
\end{fr}

\mermaidfig{p07-reshape-or-permute}{Jeden ciąg liczb, dwa kształty nad nim i
to, czego żaden z nich nie kosztuje.}{reshape albo permute}

\begin{fr}
To wystarcza, żeby postawić regułę, i warto spróbować postawić ją przed
przeczytaniem.

Rozbicie ostatniej osi na dwie jest bezpieczne. Scalenie dwóch sąsiadów jest
bezpieczne. Poproszenie o $(b, h, s, d)$ wprost z $(b, s, h \times d)$ nie jest.
Powiedz, co mają dwa bezpieczne, a czego nie ma trzecie.
\dotline
\nextframe
\end{fr}

\begin{fr}
\begin{ansblock}
Reshape jest bezpieczny dokładnie tam, gdzie rozbijane albo scalane osie
sąsiadują ze sobą i są już w kolejności, którą ma bufor.
\end{ansblock}

Wszystko, co przeniosłoby oś tam, gdzie jej nie ma, wymaga permute, a permute
musi być pierwszy. Reguła jest rozstrzygalna z samego kształtu i samego
działania, i to czyni ją wartą posiadania, zamiast pamiętania, który z dwóch
powyższych wierszy był tym właściwym.

\begin{trapbox}
\textbf{\enquote{Ten sam kształt} to nie \enquote{ta sama tablica}, a to jedyne
miejsce, w którym sygnatura typu przestaje pomagać.}

Program~\ref{prog:P06} uczynił kształt sygnaturą typu, a błąd kształtu błędem
typu, co jest mocną gwarancją i kończy się dokładnie tutaj. Oba wyrażenia wyżej
mają typ $(b, h, s, d) $. Jedno z nich to głowice twojego modelu, a drugie to
te same liczby przeplecione.

System typów, który ich nie odróżnia, nie jest zepsuty \dash{} te dwie rzeczy
naprawdę mają ten sam typ. Różnią się tym, który indeks co znaczy, a to jest
niesione w nazwach, których użyłeś, w kolejności reshape i nigdzie, gdzie
maszyna mogłaby to przeczytać.
\end{trapbox}

\begin{aibox}
Co jest całym podziałem na głowice, w zdaniu, które umiesz już sprawdzić.

Model liczy zapytania, klucze i wartości jednym szerokim iloczynem macierzy na
pozycję \dash{} warstwą Programu~\ref{prog:P06}, z osią partii z przodu.
Głowice nie są osobnym rachunkiem: szeroki wektor zostaje pocięty na $h$
plasterków po $d$, oś głowic przenosi się przed oś pozycji, a to, co było jednym
dużym iloczynem macierzy, staje się $h$ małymi wykonywanymi równolegle.

Potem $\code{'bhqd,bhkd->bhqk'}$ zwęża ostatnią oś i tylko ostatnią, więc każda
pozycja zapytania spotyka każdą pozycję klucza, osobno dla każdej głowicy i
każdego elementu partii. \textbf{Cztery osie, jedno zwężenie i żaden krok tego
nie jest nowy.} Program~\ref{prog:P32} składa resztę.
\end{aibox}
\end{fr}

\begin{fr}
Ostatnia ramka, i jest o tym, co kupiła ta notacja.

Sześć sekcji temu tablica była obrazkiem z trzema kierunkami, a kształt
księgowością. Jest krotką indeksów i regułą, działanie na niej jest
stwierdzeniem o tym, które wskaźniki się spotykają, a każda z porażek w tym
programie wzięła się z kształtu, który był legalny i nie był tym, o co
chodziło: wskaźnik wsadowy odczytany jako zwężenie, $(n)$ zestawione z
$(n, 1)$, reshape poproszony o kolejność, której bufor nie miał.

\textbf{Żadna z tych trzech nie jest błędem i być nim nie może.} Biblioteka
robi dokładnie to, co mówią kształty. Notacja wskaźnikowa jest sposobem
powiedzenia, o co ci chodziło, i dlatego ten program spędził pięć ramek na
sigmie, zanim cokolwiek wpisał.

Program~\ref{prog:P08} wraca do dwóch indeksów i zadaje o nich trudniejsze
pytanie: ile macierz naprawdę robi, co okazuje się liczbą mniejszą niż jej
kształt.
\end{fr}

% ============================================================
\begin{summarybox}
\sumitem{1--3}{\result{\textbf{Oś jest pozycją w krotce indeksów, a nie kierunkiem w
  przestrzeni.} Kształt to krotka długości, a obie konwencje partii \dash{}
  $k \times n$ i $n \times k$ \dash{} są jednym przekształceniem zapisanym na
  dwa sposoby, bo $XW\T$ jest transpozycją $WX$}.}
\sumitem{4--5}{\result{\textbf{Trzy rzeczy nazywają się wymiarem} i odpowiadają na
  różne pytania: liczba osi ($\val{p07.emb.rank}$ dla kształtu
  $(\val{p07.emb.b}, \val{p07.emb.s}, \val{p07.emb.d})$), długość jednej osi
  oraz wymiar przestrzeni, w której żyje zapisany wektor, czyli
  $\val{p07.emb.d}$, będący faktem o jednej osi, a nie o tablicy}.}
\sumitem{6--8}{\result{\emph{Rząd} tablicy to liczba jej osi, czyli inne słowo niż to z
  Programu~\ref{prog:P08}. Pod spodem jest jeden bufor
  $\val{p07.emb.total}$ liczb plus \textbf{reguła zamieniająca krotkę indeksów
  na pozycję}, i bez niej sekcji 5 nie da się postawić}.}
\sumitem{9--12}{\result{\textbf{Wskaźnik, który pojawia się dwa razy, jest sumowany, a
  ten, który pojawia się raz, przeżywa}, więc kształt wyniku odczytuje się z
  liter bez żadnej arytmetyki. $\sum_k A_{ik}B_{jk}$ to $AB\T$, a transpozycja
  nigdy nie musi się zdarzyć}.}
\sumitem{13--15}{\result{$S_{bhqk} = \sum_d Q_{bhqd}K_{bhkd}$ to wsadowy wielogłowicowy
  wynik uwagi, czytelny z reguły, zanim padnie którekolwiek z tych słów.
  \textbf{Rozstrzyga lewa strona}: wskaźnik z niej nieobecny jest sumowany, a
  obecny jest wskaźnikiem wsadowym, choćby powtarzał się wielokrotnie}.}
\sumitem{16--18}{\result{\textbf{Napis \code{einsum} to tamta suma z pominiętą sigmą.}
  \code{'ik,kj->ij'} to $\sum_k A_{ik}B_{kj}$; strzałka nazywa tych, co
  przeżywają, a wszystko na lewo od niej, czego nie ma na prawo, jest sumowane.
  Zatem \code{'ij->ji'} transponuje, \code{'ii->'} to ślad, a
  \code{'ij,j->i'} to macierz na wektorze}.}
\sumitem{21--22}{\result{\textbf{\enquote{Powtórzony znaczy sumowany} to reguła dla
  sumy, a nie dla napisu.} W \code{'bij,bjk->bik'} sumowane jest tylko $j$;
  $b$ jest wskaźnikiem wsadowym, bo jest też po prawej. Usunięcie go z prawej
  jest legalne, zwraca właściwy kształt i dodaje do siebie wyniki należące do
  różnych przykładów}.}
\sumitem{23--25}{\result{\textbf{Rozgłaszanie ma dokładną regułę}: zestaw kształty od
  prawej, uzupełnij krótszy wartościami $1$, a para zgadza się, gdy długości są
  równe albo gdy jedna wynosi $1$. $(3)$ dodane do $(3,3)$ idzie wzdłuż wierszy, a
  $(3,1)$ w dół kolumn; oba są legalne i tylko jedno jest zamierzone}.}
\sumitem{25--29}{\result{\textbf{$(n)$ przeciwko $(n,1)$ daje $(n,n)$}, więc strata
  jest uśredniana po wszystkich parach, a nadwyżka jest dokładna:
  $2\operatorname{Cov}(p,t)$. Przy doskonałym dopasowaniu raportowana strata
  wynosi $2\operatorname{Var}(t)$ i niżej nie zejdzie. Zmierzone: raportowana
  liczba jest $\val{p07.mse.ratio.hi}$ razy większa od prawdziwej przy słabym
  dopasowaniu i $\val{p07.mse.ratio.lo}$ razy przy dobrym \dash{} \textbf{im
  lepszy model, tym bardziej kłamie}}.}
\sumitem{30--32}{\result{\code{keepdims} zostawia oś długości $1$ na miejscu, żeby
  następne zestawienie było wypowiedziane, a nie wywnioskowane. A
  \textbf{żadne z reshape, transpozycji i permute nie przesuwa liczby}: każde
  oddaje ten sam bufor z inną regułą czytania, a dane przesuwa dopiero żądanie
  wyniku spójnego w pamięci}.}
\sumitem{33--36}{\result{\textbf{Ten sam kształt to nie ta sama tablica.}
  \code{x.reshape(b, s, h, d).transpose(1, 2)} i \code{x.reshape(b, h, s, d)}
  dają oba $(b, h, s, d)$ i różnią się na $\val{p07.split.differ}$ z
  $\val{p07.split.entries}$ wpisów. Reshape jest bezpieczny dokładnie tam,
  gdzie rozbijane albo scalane osie sąsiadują i są już w kolejności bufora}.}
\end{summarybox}

% ============================================================
\begin{testexercises}
\item Tablica ma kształt $(4, 6, 8)$. Podaj liczbę jej osi, liczbę wpisów i
  wymiar przestrzeni, w której żyje zapisany wektor.
  \answerto{Trzy osie; $4 \times 6 \times 8 = 192$ wpisów; a wektory mają po
  $8$ współrzędnych, żyją więc w przestrzeni o $8$ wymiarach. Trzy różne
  pytania i trzy różne liczby.}
\item Powiedz, co liczy $\sum_j A_{ij} B_{kj}$, i podaj kształt wyniku, jeśli
  $A$ jest $m \times n$, a $B$ jest $p \times n$.
  \answerto{$AB\T$, o kształcie $m \times p$. Powtórzone $j$ jest sumowane, a
  $i$ oraz $k$ przeżywają, w tej kolejności.}
\item Napisz napis \code{einsum} dla partii iloczynów macierzy, z partią z
  przodu.
  \answerto{\code{'bij,bjk->bik'}. Sumowane jest tylko $j$; $b$ jest też po
  prawej, jest więc wskaźnikiem wsadowym, a nie zwężeniem.}
\item Kształty $(5, 1, 4)$ i $(3, 4)$ są łączone po współrzędnych. Podaj
  kształt wyniku albo powiedz, dlaczego go nie ma.
  \answerto{$(5, 3, 4)$. Po uzupełnieniu do $(1, 3, 4)$ pary od prawej to
  $4$ z $4$, $1$ z $3$ i $5$ z $1$ \dash{} wszystkie trzy się zgadzają, a dwie
  z nich rozciągają.}
\item Kształty $(5, 4)$ i $(5)$ są łączone po współrzędnych. Powiedz, co się
  dzieje.
  \answerto{Odrzucone. Zestawione od prawej pary to $4$ z $5$, co nie jest ani
  równe, ani $1$. Zapisanie drugiego jako $(5, 1)$ czyni to legalnym i stosuje
  je w dół kolumn.}
\item Predykcje o kształcie $(n)$ są porównywane z wartościami docelowymi o
  kształcie $(n, 1)$. Powiedz, czym jest raportowany błąd średniokwadratowy
  wobec prawdziwego.
  \answerto{Prawdziwym plus $2\operatorname{Cov}(p, t)$, bo różnica powstaje
  dla każdej pary. Przy doskonałym dopasowaniu wynosi $2\operatorname{Var}(t)$
  i niżej nie zejdzie.}
\item Powiedz, dlaczego ten błąd jest gorszy niż błąd raportujący losową
  liczbę.
  \answerto{Bo nadwyżka jest podwojoną kowariancją, a kowariancję trening
  zwiększa. Jest łagodny, dopóki model jest zły, i ostry, gdy jest dobry,
  wygląda więc jak wypłaszczenie, a nie jak usterka.}
\item Powiedz, które z reshape, transpozycji i permute przesuwa liczby w
  pamięci i co je przesuwa.
  \answerto{Żadne, w typowym przypadku: każde zwraca ten sam bufor z inną
  regułą czytania. Kopiuje dopiero żądanie wyniku spójnego w pamięci, będące
  osobnym działaniem.}
\item Tablica o kształcie $(b, s, h \times d)$ ma stać się $(b, h, s, d)$.
  Podaj działania po kolei i powiedz, dlaczego jeden reshape nie wystarczy.
  \answerto{Rozbij ostatnią oś, co daje $(b, s, h, d)$, a potem zamień dwie
  środkowe osie. Jeden reshape do $(b, h, s, d)$ prosi o liczby w kolejności,
  której bufor nie ma, i dostaje ją, bo reshape sprawdza tylko sumę wpisów.}
\item Podaj regułę rozstrzygającą, kiedy sam reshape jest bezpieczny.
  \answerto{Gdy rozbijane albo scalane osie sąsiadują ze sobą i są już w
  kolejności, którą ma bufor. Wszystko inne wymaga wcześniejszego permute.}
\end{testexercises}

% ============================================================
\begin{furtherproblems}
\item Pokaż na wskaźnikach, że $XW\T$ i $WX$ opisują tę samą warstwę.
  \answerto{$(XW\T)_{ij} = \sum_k X_{ik} W_{jk}$ oraz
  $(WX)_{ji} = \sum_k W_{jk} X_{ik}$, czyli ta sama suma. Jedno jest więc
  transpozycją drugiego i różni je tylko kolejność osi.}
\item Podaj kształt wyniku \code{'bhqd,bhkd->bhqk'} w kategoriach czterech
  długości osi i powiedz, które długości muszą się zgadzać.
  \answerto{$(b, h, q, k)$. Długości partii i głowic muszą się zgadzać między
  argumentami, bo są wspólne, a długość cech $d$ musi się zgadzać, bo jest
  zwężana. Długości zapytań i kluczy są wolne i mogą się różnić.}
\item Powiedz, co liczy \code{'bij,bjk->ik'} i dlaczego prawie nigdy nie jest
  to, o co komukolwiek chodzi.
  \answerto{Iloczyny partii wysumowane także po partii, co daje jedną macierz.
  Jest legalne, ma kształt właściwy dla pojedynczego przykładu i dodało do
  siebie wyniki należące do różnych przykładów \dash{} jest to więc cicha
  usterka, a nie błąd.}
\item Udowodnij tożsamość
  $\operatorname{mean}_{ij}(p_i - t_j)^2 = \operatorname{mean}_i(p_i - t_i)^2
   + 2\operatorname{Cov}(p, t)$.
  \answerto{Rozwiń obie strony. Lewa to
  $\overline{p^{2}} - 2\bar p\,\bar t + \overline{t^{2}}$, bo $i$ oraz $j$
  biegną niezależnie; prawa to
  $\overline{p^{2}} - 2\overline{pt} + \overline{t^{2}}$ plus
  $2(\overline{pt} - \bar p\,\bar t)$. Zgadzają się wyraz po wyrazie.}
\item Powiedz, dlaczego zachowana oś długości $1$ jest bezpieczniejsza niż oś
  brakująca.
  \answerto{Oś długości $1$ może się tylko rozciągnąć, zestawia się więc z tym,
  co naprzeciw, a zamiar jest zapisany. Oś brakująca zostaje uzupełniona
  z przodu wartością $1$, zestawia się więc z tą osią, która akurat tam jest, a
  zamiar jest wywnioskowany z rzędu drugiego argumentu.}
\item Tablica o kształcie $(2, 3, 4)$ zostaje przekształcona do $(3, 2, 4)$.
  Powiedz, czy to to samo co zamiana dwóch pierwszych osi.
  \answerto{Nie. Reshape zachowuje kolejność bufora i przeetykietowuje ją, więc
  wpis, który był pod $(0, 1, k)$, jest teraz pod $(0, 1, k)$ czytanym inaczej;
  zamiana kładzie go pod $(1, 0, k)$. Oba dają kształt $(3, 2, 4)$ i są to
  różne tablice.}
\item Powiedz, co musiałoby być prawdą o systemie typów, żeby wyłapał zły
  podział na głowice.
  \answerto{Musiałby nieść to, co każda oś \emph{znaczy}, a nie tylko jak jest
  długa \dash{} kształt $(b, h, s, d)$ z czterema nazwami, a nie z czterema
  liczbami. Oba wyrażenia mają ten sam kształt, więc żaden system rozumujący o
  samych kształtach ich nie rozdzieli.}
\item Program~\ref{prog:P03} wycenia przebieg przez pamięć. Powiedz, które z
  działań sekcji 5 go pociąga i kiedy.
  \answerto{Tylko żądanie wyniku spójnego w pamięci, i tylko wtedy, gdy bieżąca
  reguła czytania bufora nie pasuje już do kształtu. Reshape tablicy, która jest
  już spójna, nie kosztuje nic; ten sam reshape po permute musi kopiować, bo
  kolejność, której potrzebuje, nie jest tą, która tam jest.}
\end{furtherproblems}
